\subsection{\texorpdfstring{Profundidad y propiedad $(S_k)$.}{Profundidad y propiedad (S\string_k)}}

\begin{definition}[5.7.1]
\label{IV.2.s5.7.1-fr}
Sean $X$ un preesquema localmente noetheriano y $\mathscr F$ un
$\mathcal O_X$-módulo coherente. Se llama profundidad \textup{(resp.}
coprofundidad\textup{)} de $\mathscr F$ en un punto $x\in X$ al número
$\operatorname{prof}(\mathscr F_x)$ \textup{(resp.}
$\operatorname{coprof}(\mathscr F_x)$\textup{)} \textup{(\textbf{0}, 16.4.5 y 16.4.9)}. Se llama
coprofundidad de $\mathscr F$ al número
\begin{equation}
\operatorname{coprof}(\mathscr F)
=\sup_{x\in X}\operatorname{coprof}(\mathscr F_x).
\tag{7.5.1.1}\label{IV.2.s5.7.1.1-fr}
\end{equation}

Se dice que $\mathscr F$ es un $\mathcal O_X$-módulo de Cohen-Macaulay en el
punto $x$ si $\mathscr F_x$ es un $\mathcal O_x$-módulo de Cohen-Macaulay,
es decir \textup{(\textbf{0}, 16.5.1)}, si $\operatorname{coprof}(\mathscr F_x)=0$. Se dice que
$\mathscr F$ es un $\mathcal O_X$-módulo de Cohen-Macaulay si lo es en todo
punto, dicho de otro modo, si $\operatorname{coprof}(\mathscr F)=0$. Un punto
$x\in X$ tal que $\mathcal O_x$ sea un anillo de Cohen-Macaulay se llama
también punto de Cohen-Macaulay de $X$.
\end{definition}

Se llama \emph{coprofundidad} de $X$ y se denota por
$\operatorname{coprof}(X)$ el número $\operatorname{coprof}(\mathcal O_X)$.
Se dice que $X$ es un \emph{preesquema de Cohen-Macaulay} si $\mathcal O_X$
es un $\mathcal O_X$-módulo de Cohen-Macaulay, dicho de otro modo, si
$\operatorname{coprof}(X)=0$. Todo preesquema localmente noetheriano de
\emph{dimensión} 0 es evidentemente un preesquema de Cohen-Macaulay. Decir que
$\operatorname{Spec}(A)$ es un esquema de Cohen-Macaulay significa que $A$
es un anillo de Cohen-Macaulay (\textbf 0, 16.5.13).

\begin{definition}[5.7.2]
\label{IV.2.s5.7.2-fr}
Sean $X$ un preesquema localmente noetheriano, $\mathscr F$ un
$\mathcal O_X$-módulo coherente y $k$ un entero positivo o negativo. Se dice
que $\mathscr F$ posee la propiedad $(S_k)$ si, para todo $x\in X$, se tiene
\begin{equation}
\operatorname{prof}(\mathscr F_x)\geq
\inf(k,\dim(\mathscr F_x))
\tag{5.7.2.1}\label{IV.2.s5.7.2.1-fr}
\end{equation}
Se dice que $\mathscr F$ posee la propiedad $(S_k)$ en un punto $x\in X$ si,
para toda generalización $x'$ de $x$ en $X$, se tiene
\begin{equation}
\operatorname{prof}(\mathscr F_{x'})\geq
\inf(k,\dim(\mathscr F_{x'})).
\tag{5.7.2.2}\label{IV.2.s5.7.2.2-fr}
\end{equation}
\end{definition}

Se dice que $X$ verifica la propiedad $(S_k)$ (resp. verifica la propiedad
$(S_k)$ en un punto $x$) si $\mathcal O_X$ verifica la propiedad $(S_k)$
(resp. verifica la propiedad $(S_k)$ en el punto $x$).

Para que $\mathscr F$ verifique la propiedad $(S_k)$, es evidentemente
necesario y suficiente que la verifique en todo punto de $X$. Si $U$ es un
abierto de $X$ y $\mathscr F$ verifica $(S_k)$, lo mismo ocurre con
$\mathscr F|U$; recíprocamente, si $(U_\alpha)$ es un recubrimiento abierto
de $X$ y si $\mathscr F|U_\alpha$ verifica $(S_k)$ para todo $\alpha$,
$\mathscr F$ verifica $(S_k)$.
